\chapter{2005 Nobel Prize in Economics}
\textbf{Laureates: }Robert Aumann (Israel), Thomas Schelling (USA)

\section*{Reason for Award}
For having enhanced our understanding of conflict and cooperation through
game-theory analysis

\section{What They Did}
Aumann and Schelling applied game theory to analyze conflict and cooperation,
transforming it from an abstract mathematical tool into a practical framework
for understanding strategic interaction across a wide range of domains.

Aumann, born in Germany in 1930, made foundational contributions to the
mathematical theory of repeated games. In his 1959 paper, "Repeated Games
with Perfect Monitoring," he showed that in infinitely repeated prisoner's
dilemmas, cooperative outcomes can be sustained as equilibrium strategies when
players are sufficiently patient. The key insight is that the threat of future
punishment for current defection can deter short-term deviation. This result
shows how cooperation can emerge in situations where one-shot game theory
predicts defection.

Aumann also developed the concept of common knowledge, formalizing the
conditions under which players can coordinate on equilibria in coordination
games. Common knowledge requires not only that all players know a fact, but
that they know that everyone knows it, and know that everyone knows that
everyone knows it, ad infinitum. This infinite hierarchy of knowledge is
essential for coordinating expectations and achieving equilibrium outcomes
in games with multiple equilibria.

His introduction of correlated equilibrium in 1974 generalized the Nash
equilibrium concept by allowing players to coordinate their strategies
through signals from a common source, such as a mediator or public signal.
Unlike Nash equilibrium, where players choose strategies independently,
correlated equilibrium allows for correlated recommendations that players
have incentive to follow. This concept has applications to auction design,
traffic management, and market coordination.

Schelling, born in Colorado in 1921, applied game-theoretic insights to
military strategy, bargaining, and conflict resolution. In his seminal work
"The Strategy of Conflict" (1960), he analyzed how commitment mechanisms
shape bargaining outcomes. He showed that the ability to constrain one's
own future choices--through irreversible investments, public announcements,
or institutional arrangements--can strengthen one's negotiating position
by making threats credible.

Schelling's insights about focal points showed how players coordinate in
the absence of communication. Focal points are salient solutions that
players naturally converge on without explicit coordination. His analysis
of credibility, commitment, and brinkmanship informed understanding of
nuclear deterrence, international negotiations, and labor disputes.

\section{Impact on Economic Thinking}
Aumann and Schelling's contributions transformed game theory from an abstract
mathematical discipline into a practical tool for analyzing real-world
strategic interaction. Aumann's mathematical rigor provided the formal
foundation, while Schelling's intuitive insights demonstrated the relevance
of game-theoretic reasoning to politics, diplomacy, and everyday social situations.

Repeated game theory showed how cooperation can be sustained in ongoing
relationships where reputation matters. Schelling's insights about commitment
revealed how the ability to bind oneself strategically can improve outcomes
for all parties. These insights apply to business relationships, international
negotiations, labor negotiations, and family dynamics.

The correlated equilibrium concept enriched strategic analysis, allowing for
coordination through communication and mediation without requiring binding
contracts. Common knowledge formalized the coordination requirement for shared
understanding, with applications ranging from focal points in coordination
games to the analysis of self-fulfilling prophecies and financial runs.

Their work extended beyond economics into political science, international
relations, evolutionary biology, and philosophy. Schelling's work on nuclear
deterrence during the Cold War had real-world policy implications that
extended well beyond academic economics. His insights about credible
commitment and brinkmanship informed US strategy during the Cold War.

Aumann's work on epistemic game theory provided rigorous foundations for
analyzing knowledge and belief in strategic situations. His contributions
to the theory of learning in games have influenced evolutionary game theory
and the study of adaptive behavior.

\section{Modern Perspectives Today}
Repeated game theory extends to evolutionary biology, where replicator dynamics
shows how cooperative strategies can evolve under natural selection. Behavioral
experiments validate Aumann's cooperation findings in laboratory settings,
while also revealing conditions under which cooperation breaks down.

Correlated equilibrium now facilitates coordination in traffic management,
spectrum allocation, and auction design. Mechanism design incorporates
correlated signals to improve allocation efficiency. Matching markets for
school assignments, residency placements, and kidney transplants draw on
these insights to design stable and efficient allocation procedures.

The framework informs industrial organization, labor negotiations, and
corporate governance. Coalition formation, alliances, joint ventures, and
partnerships all employ game-theoretic analysis. The study of international
institutions, treaty compliance, and enforcement mechanisms builds on Aumann
and Schelling's insights about repeated interaction and credible commitment.

Their work has influenced the design of institutions for managing common
pool resources, international environmental agreements, and global governance.
The insights about commitment and coordination remain relevant for addressing
collective action problems in an increasingly interconnected world.
